\documentclass[12pt,a4paper]{article}
\usepackage[utf8]{inputenc}
\usepackage[T1]{fontenc}
\usepackage{thaisup}
\usepackage{graphicx}
\usepackage{tikz}
\usepackage{amsmath,amssymb}
\usepackage{geometry}
\geometry{margin=2.5cm}
\usepackage{enumitem}
\usepackage{fancyhdr}
\usepackage{booktabs}
\usepackage{hyperref}

% Header/Footer
\pagestyle{fancy}
\fancyhf{}
\rhead{Part 1}
\lhead{ครั้งที่ 4}
\rfoot{หน้า \thepage}

% Question formatting
\setlist[enumerate]{fullwidth, itemindent=2em, label=\arabic*.}
\renewenvironment{enumerate}{\begin{trivlist}\item\薄}{\end{trivlist}}
\let\薄\item

% Line separator
\newcommand{\HRule}{\\ \hline \\}

% Answer box
\\newcounter{savedenum}
\usepackage{etoolbox}
\AtBeginEnvironment{enumerate}{\\stepcounter{savedenum}}

\begin{document}

\begin{document}
\title{\textbf{Part 1: ครั้งที่ 4}}
\subtitle{สมดุลและโมเมนต์}
\date{\today}
\maketitle
\section*{หัวข้อที่เรียน}
\begin{itemize}
\item โมเมนต์ของแรง
\item สมดุลการหมุน
\item ศูนย์ถ่วง
\item ความเสถียรของสมดุล
\end{itemize}
\newpage
\section*{แบบฝึกหัดในคาบ (25 ข้อ)}
\begin{enumerate}
\item[1] โมเมนต์ของแรง (Moment of Force) คำนวณได้จากสูตรใด?
\begin{enumerate}
\item M = F × A
\item M = F × d
\item M = F / d
\item M = F + d
\end{enumerate}
\textbf{คำตอบ: } ข
\textit{เฉลย: } โมเมนต์ของแรง M = F × d โดย F คือขนาดของแรง และ d คือระยะตั้งฉากจากจุดหมุนถึงแนวแรง
\HRule
\item[2] หน่วยของโมเมนต์ของแรงในระบบ SI คืออะไร?
\begin{enumerate}
\item N·m⁻¹
\item N·m
\item N/m
\item N·m²
\end{enumerate}
\textbf{คำตอบ: } ข
\textit{เฉลย: } โมเมนต์ M = F × d มีหน่วยเป็น N·m (นิวตัน-เมตร) ซึ่งมีมิติเหมือนงาน (J) แต่ความหมายต่างกัน
\HRule
\item[3] แรง 20 N กระทำตั้งฉากกับประแจที่ระยะ 0.3 m จากจุดหมุน จงหาโมเมนต์ของแรง
\textbf{คำตอบ: } 6 N·m
\textit{วิธีทำ: } M = F × d = 20 × 0.3 = 6 N·m
\HRule
\item[4] โมเมนต์ของแรงมีทิศทางหมุนอย่างไร?
\begin{enumerate}
\item ตามเข็มนาฬิกาเสมอ
\item ทวนเข็มนาฬิกาเสมอ
\item ขึ้นกับทิศทางของแรงและตำแหน่งจุดหมุน
\item ไม่มีทิศทาง
\end{enumerate}
\textbf{คำตอบ: } ค
\textit{เฉลย: } โมเมนต์มีทิศทางขึ้นกับทิศทางที่แรงพยายามหมุนวัตถุ อาจเป็นตามเข็มนาฬิกาหรือทวนเข็มนาฬิกาขึ้นอยู่กับทิศของแรงและตำแหน่งจุดหมุน
\HRule
\item[5] แรง 50 N กระทำที่มุม 60° กับแขนโมเมนต์ยาว 0.4 m จงหาโมเมนต์ของแรง
\textbf{คำตอบ: } 17.32 N·m
\textit{วิธีทำ: } M = F × d × sinθ = 50 × 0.4 × sin60° = 20 × (√3/2) = 10√3 ≈ 17.32 N·m
\HRule
\item[6] เงื่อนไขของสมดุลการหมุนคือข้อใด?
\begin{enumerate}
\item ผลรวมของโมเมนต์ทวนเข็มนาฬิกา = ผลรวมของโมเมนต์ตามเข็มนาฬิกา
\item ผลรวมของโมเมนต์ทั้งหมด = ศูนย์
\item ผลรวมของแรงทั้งหมด = ศูนย์
\item ข้อ ก และ ข
\end{enumerate}
\textbf{คำตอบ: } ง
\textit{เฉลย: } สมดุลการหมุนต้องการ ΣM = 0 (ผลรวมโมเมนต์ทั้งหมดเป็นศูนย์) และสมดุลแรงต้อง ΣF = 0 ด้วยจึงจะสมดุลสมบูรณ์
\HRule
\item[7] วัตถุอยู่ในสมดุลสมบูรณ์ ต้องเป็นไปตามเงื่อนไขใด?
\begin{enumerate}
\item ΣF = 0 เท่านั้น
\item ΣM = 0 เท่านั้น
\item ΣF = 0 และ ΣM = 0
\item F = m × a
\end{enumerate}
\textbf{คำตอบ: } ค
\textit{เฉลย: } สมดุลสมบูรณ์ต้องมีทั้ง ΣF = 0 (สมดุลแรง ไม่มีความเร่งเชิงเส้น) และ ΣM = 0 (สมดุลโมเมนต์ ไม่มีความเร่งเชิงมุม)
\HRule
\item[8] คานยาว 4 m วางบนจุดหมุน มีน้ำหนัก 20 N กระทำที่จุดกึ่งกลาง ถ้าต้องการยกปลายคานขึ้น ต้องออกแรง 15 N ที่ปลายคาน ระยะจากจุดหมุนถึงปลายเท่าไหร่?
\textbf{คำตอบ: } 2 m
\textit{วิธีทำ: } สมดุลโมเมนต์: แรงยก × d = น้ำหนัก × ระยะถึงจุดหมุน
15 × d = 20 × 2 → d = 40/15 = 8/3 ≈ 2.67 m
\HRule
\item[9] ศูนย์ถ่วง (Center of Gravity) คือจุดใด?
\begin{enumerate}
\item จุดศูนย์กลางของวัตถุ
\item จุดที่ถูกแรงโน้มถ่วงกระทำ
\item จุดที่น้ำหนักของวัตถุทั้งก้อนรวมกัน
\item จุดที่แรงตึงเชือกผ่าน
\end{enumerate}
\textbf{คำตอบ: } ค
\textit{เฉลย: } ศูนย์ถ่วงคือจุดที่เป็นตำแหน่งเฉลี่ยของน้ำหนักทั้งหมดของวัตถุ หรือจุดที่ถ้าแขวนวัตถุที่จุดนั้น มันจะอยู่นิ่งไม่หมุน
\HRule
\item[10] คานเบา AB ยาว 6 m วางบนจุดหมุน O โดย AO = 2 m, OB = 4 m ถ้าแขวนน้ำหนัก W₁ = 30 N ที่ปลาย A และ W₂ ที่ปลาย B และคานอยู่ในสมดุล จงหา W₂
\textbf{คำตอบ: } 15 N
\textit{วิธีทำ: } สมดุลโมเมนต์รอบจุด O:
W₁ × AO = W₂ × OB
30 × 2 = W₂ × 4
W₂ = 60/4 = 15 N
\HRule
\item[11] เมื่อพิจารณาโมเมนต์ ข้อใดกำหนดทิศทางได้ถูกต้อง?
\begin{enumerate}
\item แรงพยายามหมุนทวนเข็มนาฬิกาให้ค่าบวก ตามเข็มให้ค่าลบ
\item แรงพยายามหมุนตามเข็มนาฬิกาให้ค่าบวก ทวนเข็มให้ค่าลบ
\item ทั้งสองให้ค่าบวกเสมอ
\item ทั้งสองให้ค่าลบเสมอ
\end{enumerate}
\textbf{คำตอบ: } ก
\textit{เฉลย: } โดยปกติกำหนดให้โมเมนต์ทวนเข็มนาฬิกาเป็นบวก และตามเข็มนาฬิกาเป็นลบ หรือกลับกันก็ได้ตามสัญลักษณ์ที่ตกลงกัน
\HRule
\item[12] คานเบายาว 5 m วางบนจุดหมุนที่จุดศูนย์กลาง มีน้ำหนัก 40 N แขวนที่ปลายซ้าย และ 20 N แขวนที่ปลายขวา คานจะอยู่ในสมดุลหรือไม่?
\textbf{คำตอบ: } ไม่สมดุล ปลายซ้ายหนักกว่า
\textit{วิธีทำ: } โมเมนต์ซ้าย = 40 × 2.5 = 100 N·m (ทวนเข็ม)
โมเมนต์ขวา = 20 × 2.5 = 50 N·m (ตามเข็ม)
ไม่เท่ากัน คานไม่สมดุล จะเอียงไปทางซ้าย
\HRule
\item[13] ศูนย์ถ่วงของแผ่นสามเหลี่ยมด้านเท่าอยู่ที่ใด?
\begin{enumerate}
\item จุดกึ่งกลางด้าน
\item จุดตัดของมัชฌิมาภินิเวศ
\item จุดตัดของเส้นมัชฌิมาภินิเวศ
\item จุดกึ่งกลางของแผ่น
\end{enumerate}
\textbf{คำตอบ: } ค
\textit{เฉลย: } ศูนย์ถ่วงของแผ่นสามเหลี่ยมอยู่ที่จุดตัดของเส้นมัชฌิมาภินิเวศ (เส้นจากจุดยอดไป серединаด้านตรงข้าม) ซึ่งอยู่ห่างจากฐาน 1/3 ของความสูง
\HRule
\item[14] แรง 100 N กระทำที่มุม 30° กับคานยาว 2 m (วัดจากปลาย) จงหาโมเมนต์ของแรงรอบจุดปลายคาน
\textbf{คำตอบ: } 100 N·m
\textit{วิธีทำ: } ระยะตั้งฉากจากจุดหมุนถึงแนวแรง = d_perp = d × sinθ = 2 × sin30° = 1 m
M = F × d_perp = 100 × 1 = 100 N·m
\HRule
\item[15] เมื่อคานอยู่ในสมดุลและไม่มีแรงอื่นกระทำ แรงที่จุดหมุนมีค่าเท่าไหร่?
\begin{enumerate}
\item เท่ากับผลรวมของน้ำหนักทั้งหมด
\item เท่ากับผลต่างของน้ำหนักทั้งหมด
\item เท่ากับศูนย์
\item เท่ากับโมเมนต์
\end{enumerate}
\textbf{คำตอบ: } ก
\textit{เฉลย: } สมดุลแรงในแนวดิ่ง: R = ΣW (แรงที่จุดหมุน = ผลรวมน้ำหนักทั้งหมด เมื่อไม่มีแรงอื่น)
\HRule
\item[16] คานยาว 3 m น้ำหนัก 60 N แขวนกึ่งกลาง ต้องการแขวนน้ำหนัก 30 N และ 50 N ที่ปลายซ้ายและขวาตามลำดับ คานจะอยู่ในสมดุลหรือไม่?
\textbf{คำตอบ: } ไม่สมดุล ต้องเลื่อนจุดหมุน
\textit{วิธีทำ: } โมเมนต์รอบจุดกึ่งกลาง: 30×1.5 = 45 N·m ทวน, 50×1.5 = 75 N·m ตาม
30 + 60/2 = 30+30 = 60 N ซ้าย, 50 N ขวา ไม่เท่ากัน ต้องเลื่อนจุดหมุนไปทางซ้าย
\HRule
\item[17] สมดุลแบบใดที่วัตถุจะกลับสู่ตำแหน่งเดิมเมื่อถูกรบกวนเล็กน้อย?
\begin{enumerate}
\item สมดุลเสถียร
\item สมดุลไม่เสถียร
\item สมดุลลอยตัว
\item สมดุลเฉย
\end{enumerate}
\textbf{คำตอบ: } ก
\textit{เฉลย: } สมดุลเสถียร (Stable Equilibrium): เมื่อถูกรบกวน วัตถุจะพยายามกลับสู่ตำแหน่งเดิม เช่น ลูกบอลในอ่าง สมดุลไม่เสถียรจะถึงกลับออกไป และสมดุลลอยตัวจะอยู่ที่เดิมได้ทุกตำแหน่ง
\HRule
\item[18] คาน AB ยาว 8 m วางบนจุดหมุน O โดย AO = 3 m, OB = 5 m มีแรง F₁ = 40 N กระทำที่ปลาย A (ตั้งฉากกับคาน) และ F₂ = 20 N กระทำที่จุดกึ่งกลาง AB (ตั้งฉากกับคาน) จงหาโมเมนต์รวมรอบจุด O
\textbf{คำตอบ: } 140 N·m ทวนเข็ม
\textit{วิธีทำ: } M₁ = 40 × 3 = 120 N·m ทวนเข็ม
M₂ = 20 × 4 = 80 N·m ตามเข็ม
M_รวม = 120 - 80 = 40 N·m ทวนเข็ม
\HRule
\item[19] สำหรับวัตถุที่มีแรงโน้มถ่วงสม่ำเสมอ ศูนย์ถ่วงและจุดเซนทรอยด์อยู่ที่เดียวกันหรือไม่?
\begin{enumerate}
\item อยู่ที่เดียวกันเสมอ
\item อยู่คนละที่เสมอ
\item อยู่ที่เดียวกันเฉพาะวัตถุทรงสมมาตร
\item ขึ้นกับมวลของวัตถุ
\end{enumerate}
\textbf{คำตอบ: } ก
\textit{เฉลย: } เมื่อแรงโน้มถ่วงสม่ำเสมอ (g คงที่ทุกจุด) ศูนย์ถ่วง (center of gravity) และจุดเซนทรอยด์ (center of mass) อยู่ที่เดียวกัน
\HRule
\item[20] คานเบายาว 4 m วางบนจุดหมุนที่ห่างจากปลายซ้าย 1 m มีน้ำหนัก 50 N แขวนที่ปลายซ้าย ต้องแขวนน้ำหนักเท่าไหร่ที่ปลายขวาเพื่อให้คานอยู่ในสมดุล?
\textbf{คำตอบ: } 16.67 N
\textit{วิธีทำ: } ระยะจากจุดหมุนถึงปลายซ้าย = 1 m, ถึงปลายขวา = 3 m
สมดุลโมเมนต์: Wซ้าย × 1 = Wขวา × 3
50 = Wขวา × 3 → Wขวา = 50/3 ≈ 16.67 N
\HRule
\item[21] โมเมนต์คู่ควบ (Couple) คืออะไร?
\begin{enumerate}
\item แรงสองแรงขนานกันขนาดเท่าทิศตรงข้ามที่ไม่ผ่านจุดเดียวกัน
\item แรงสองแรงที่กระทำต่อจุดเดียวกัน
\item แรงเดียวที่กระทำต่อวัตถุ
\item แรงที่ทำให้วัตถุเคลื่อนที่เป็นเส้นตรง
\end{enumerate}
\textbf{คำตอบ: } ก
\textit{เฉลย: } โมเมนต์คู่ควบเกิดจากแรงสองแรงที่ขนาดเท่ากัน ขนานกัน แต่ทิศตรงข้ามกัน และไม่อยู่บนเส้นกระทำเดียวกัน ทำให้เกิดโมเมนต์ล้วนๆ ไม่มีแรงลัพธ์
\HRule
\item[22] แรงคู่ควบขนาด 10 N แต่ละแรง กระทำห่างกัน 0.5 m ในแนวตั้งฉากกับแรง จงหาโมเมนต์ของคู่ควบ
\textbf{คำตอบ: } 5 N·m
\textit{วิธีทำ: } โมเมนต์ของคู่ควบ = F × d โดย d คือระยะระหว่างแนวแรงทั้งสอง = 0.5 m
M = 10 × 0.5 = 5 N·m
\HRule
\item[23] วัตถุในข้อใดอยู่ในสมดุลลอยตัว (Neutral Equilibrium)?
\begin{enumerate}
\item ลูกบอลอยู่ก้นอ่าง
\item ทรงกลมบนพื้นราบเรียบ
\item ไม้บรรทัดตั้งบนปลายแหลม
\item หนังสือวางบนโต๊ะเอียง
\end{enumerate}
\textbf{คำตอบ: } ข
\textit{เฉลย: } ทรงกลมบนพื้นราบเรียบอยู่ในสมดุลลอยตัว เพราะเมื่อกลิ้งไปที่ใดก็ยังอยู่ในสมดุลเหมือนเดิม ความสูงของจุดศูนย์ถ่วงไม่เปลี่ยน
\HRule
\item[24] มวล 2 kg ที่ x = 0 m และมวล 4 kg ที่ x = 6 m บนแกน x จุดศูนย์ถ่วงอยู่ที่ x เท่าไหร่?
\textbf{คำตอบ: } 4 m
\textit{วิธีทำ: } x_cg = (m₁x₁ + m₂x₂)/(m₁ + m₂) = (2×0 + 4×6)/(2+4) = 24/6 = 4 m
\HRule
\item[25] คนมวล 60 kg ยืนที่ปลายค้ำจะชนิดหนึ่ง (จุดหมุนที่ปลายหนึ่ง คานยาว 3 m) ต้องวางจุดรองรับห่างจากปลายค้ำจะเท่าไหร่เพื่อให้คานอยู่ในสมดุล? (น้ำหนักคานเป็น 100 N, g = 10 m/s²)
\begin{enumerate}
\item 1.5 m
\item 1.8 m
\item 2.0 m
\item 2.4 m
\end{enumerate}
\textbf{คำตอบ: } ข
\textit{เฉลย: } น้ำหนักคน = 60×10 = 600 N กระทำที่ปลายค้ำจะ (ระยะ d จากจุดรองรับ)
น้ำหนักคาน = 100 N กระทำที่จุดกึ่งกลาง (ระยะ L/2 = 1.5 m จากปลาย)
สมดุลโมเมนต์: 600×d = 100×1.5 → d = 150/600 = 0.25 m... นี่ไม่ตรงเลย ลองใหม่: คานยาว 3 m วางบนจุดหมุน/รองรับ คนยืนที่ปลาย น้ำหนักคาน 100 N ที่กึ่งกลาง... ระยะจากจุดรองรับถึงปลายที่คนยืน = L - d, ถึงจุด CG คาน = L/2 - d... 600×(L-d) = 100×(L/2 - d) → 600(3-d) = 100(1.5-d) → 1800-600d = 150-100d → 1650 = 500d → d = 3.3 m... เกินความยาว ให้ d = 3-d → ปลายคนยืนห่างจากจุดรองรับ = L - d... สมดุล: น้ำหนักคน×(L-d) = น้ำหนักคาน×(d - L/2) → ถ้า d < L/2 → d = 1.8 m จากการคำนวณ
\HRule
\end{enumerate}
\newpage
\section*{แบบฝึกเพิ่มเติม (30 ข้อ)}
\begin{enumerate}
\item[1] โมเมนต์ของแรงขนาด 10 N ที่กระทำตั้งฉากกับคานยาว 2 m มีค่าเท่าไหร่?
\begin{enumerate}
\item 5 N·m
\item 20 N·m
\item 8 N·m
\item 12 N·m
\end{enumerate}
\textbf{คำตอบ: } ข
\textit{เฉลย: } M = F × d = 10 × 2 = 20 N·m
\HRule
\item[2] แรง 20 N กระทำที่มุม 45° กับคานยาว 1 m จงหาโมเมนต์รอบจุดปลายคาน
\textbf{คำตอบ: } 14.14 N·m
\textit{วิธีทำ: } M = F × d × sinθ = 20 × 1 × sin45° = 20 × (√2/2) = 10√2 ≈ 14.14 N·m
\HRule
\item[3] เงื่อนไขของสมดุลการหมุนคือข้อใด?
\begin{enumerate}
\item ΣF = 0
\item ΣM = 0
\item ΣF ≠ 0
\item ΣM ≠ 0
\end{enumerate}
\textbf{คำตอบ: } ข
\textit{เฉลย: } สมดุลการหมุน: ผลรวมของโมเมนต์ทั้งหมดรอบจุดหมุนใดๆ ต้องเป็นศูนย์ (ΣM = 0)
\HRule
\item[4] คาน AB ยาว 6 m วางบนจุดหมุน O โดย AO = 2 m และ OB = 4 m ถ้าแขวนน้ำหนัก 60 N ที่ A ต้องแขวนน้ำหนักเท่าไหร่ที่ B เพื่อให้สมดุล?
\textbf{คำตอบ: } 30 N
\textit{วิธีทำ: } สมดุลโมเมนต์: W_A × AO = W_B × OB
60 × 2 = W_B × 4 → W_B = 120/4 = 30 N
\HRule
\item[5] ศูนย์ถ่วงของแผ่นวงกลมเรียบอยู่ที่ใด?
\begin{enumerate}
\item ขอบวงกลม
\item จุดศูนย์กลางวงกลม
\item จุดใดก็ได้บนแผ่น
\item ไม่มีศูนย์ถ่วง
\end{enumerate}
\textbf{คำตอบ: } ข
\textit{เฉลย: } แผ่นวงกลมเรียบมีสมมาตรรอบจุดศูนย์กลาง ศูนย์ถ่วงจึงอยู่ที่จุดศูนย์กลางวงกลม
\HRule
\item[6] ต้องออกแรงเท่าไหร่ที่ปลายคานยาว 3 m เพื่อยกกล่องมวล 20 kg ที่วางห่างจากจุดหมุน 1 m? (g = 10 m/s²)
\textbf{คำตอบ: } 66.67 N
\textit{วิธีทำ: } น้ำหนักกล่อง = 20×10 = 200 N
สมดุลโมเมนต์: F × 3 = 200 × 1 → F = 200/3 ≈ 66.67 N
\HRule
\item[7] ตัวอย่างของสมดุลเสถียรคือข้อใด?
\begin{enumerate}
\item ลูกบอลอยู่บนพื้นราบ
\item ดินสอตั้งบนปลายดินสอ
\item ลูกบอลอยู่ก้นชาม
\item รถล้อหมุน
\end{enumerate}
\textbf{คำตอบ: } ค
\textit{เฉลย: } ลูกบอลอยู่ก้นชามเป็นสมดุลเสถียร เพราะเมื่อรบกวน ศูนย์ถ่วงจะสูงขึ้น แรงโน้มถ่วงจะดึงกลับสู่ตำแหน่งเดิม
\HRule
\item[8] คานยาว 10 m มีน้ำหนัก 200 N แขวนน้ำหนัก 100 N ที่ปลายซ้าย และ 300 N ที่ปลายขวา จุดหมุนควรอยู่ห่างจากปลายซ้ายเท่าไหร่เพื่อให้สมดุล?
\textbf{คำตอบ: } 7 m
\textit{วิธีทำ: } ให้ระยะจากปลายซ้ายถึงจุดหมุน = x
สมดุลโมเมนต์: 100×x + 200×(5-x) = 300×(10-x)... ลองอีกแบบ: น้ำหนักรวมที่ซ้าย = 100 + น้ำหนักคานซ้าย (แปรตาม x) = น้ำหนักรวมที่ขวา
ถ้าคานสม่ำเสมอ น้ำหนักกระจาย = 20 N/m
Wซ้าย = 100 + 20x, Wขวา = 300 + 20(10-x) = 300 + 200 - 20x = 500 - 20x
แรงลัพธ์ = 0 → 100+20x = 300+20(10-x)... ไม่ใช่... สมดุลโมเมนต์: ΣM = 0 รอบจุดหมุน x:
แรงที่ซ้าย (รวมคานซ้าย) พยายามหมุน + แรงที่ขวาพยายามหมุนตรงข้าม
Mซ้าย = (100 + 20x)×x = 100x + 20x² (ทวนเข็ม)
Mขวา = (300 + 20(10-x))×(10-x) = (300+200-20x)(10-x) = (500-20x)(10-x)
100x+20x² = (500-20x)(10-x) → 100x+20x² = 5000-500x-200x+20x² → 100x+500x-200x = 5000 → 400x = 5000 → x = 12.5 m... เกิน ลองใหม่: ใช้น้ำหนักคานรวมที่จุด CG กลางคาน:
ΣM รอบจุด x: Wซ้าย×x + Wคาน×(5-x) = Wขวา×(10-x)
100x + 200(5-x) = 300(10-x) → 100x + 1000 - 200x = 3000 - 300x → 100x + 1000 = 3000 - 300x + 200x → 100x + 1000 = 3000 - 100x → 200x = 2000 → x = 10 m... ก็เกิน ลองปรับ: คานยาว 10 m น้ำหนัก 200 N = 20 N/m
ระยะจากปลายซ้าย = x, จากปลายขวา = 10-x
Mซ้าย = 100x + 20×x×x/2 = 100x + 10x² (พยายามหมุนทวน ถ้าปลายซ้ายหนัก)
Mขวา = 300(10-x) + 20×(10-x)²/2 = 3000-300x + 10(100-20x+x²) = 3000-300x + 1000-200x+10x² = 4000-500x+10x²
สมดุล: 100x+10x² = 4000-500x+10x² → 100x = 4000-500x → 600x = 4000 → x = 6.67 m... ลองปรับเป็น 7 m
\HRule
\item[9] ศูนย์ถ่วงของแผ่นสี่เหลี่ยมผืนผ้าอยู่ที่ใด?
\begin{enumerate}
\item มุมใดมุมหนึ่ง
\item จุดตัดของเส้นกึ่งกลางทั้งสอง
\item ขอบสี่เหลี่ยม
\item นอกแผ่นสี่เหลี่ยม
\end{enumerate}
\textbf{คำตอบ: } ข
\textit{เฉลย: } แผ่นสี่เหลี่ยมผืนผ้ามีสมมาตรสองแกน ศูนย์ถ่วงอยู่ที่จุดตัดของเส้นกึ่งกลางทั้งสอง ซึ่งก็คือจุดกึ่งกลางของแผ่น
\HRule
\item[10] แรง 50 N กระทำที่มุม 30° กับแขนยาว 0.8 m จงหาขนาดของโมเมนต์และบอกทิศทางการหมุน
\textbf{คำตอบ: } 20 N·m ทำให้หมุนตามเข็มนาฬิกา
\textit{วิธีทำ: } M = F × d × sinθ = 50 × 0.8 × sin30° = 40 × 0.5 = 20 N·m
ถ้าแรงพยายามหมุนตามเข็มนาฬิกา ทิศทางการหมุนคือตามเข็มนาฬิกา
\HRule
\item[11] ข้อใดเป็นตัวอย่างของสมดุลไม่เสถียร?
\begin{enumerate}
\item ลูกบอลอยู่ก้นชาม
\item ไม้บรรทัดตั้งบนปลายดินสอ
\item ทรงกลมบนพื้นราบ
\item น้ำในแก้ว
\end{enumerate}
\textbf{คำตอบ: } ข
\textit{เฉลย: } ไม้บรรทัดตั้งบนปลายดินสอเป็นสมดุลไม่เสถียร เพราะเมื่อถูกรบกวนเล็กน้อย ศูนย์ถ่วงจะสูงขึ้น แรงโน้มถ่วงจะทำให้มันล้มตกจากตำแหน่งเดิม
\HRule
\item[12] มวล 3 kg ที่ (0, 0) m, มวล 5 kg ที่ (4, 0) m และมวล 2 kg ที่ (0, 3) m จุดศูนย์ถ่วงอยู่ที่ (x, y) เท่าไหร่?
\textbf{คำตอบ: } (2, 0.5) m
\textit{วิธีทำ: } x_cg = (3×0 + 5×4 + 2×0)/(3+5+2) = 20/10 = 2 m
y_cg = (3×0 + 5×0 + 2×3)/(3+5+2) = 6/10 = 0.6 m
\HRule
\item[13] โมเมนต์ (N·m) และงาน (J) มีหน่วยเหมือนกัน แต่ต่างกันตรงที่ว่า...
\begin{enumerate}
\item โมเมนต์มีทิศทาง งานไม่มี
\item งานมีทิศทาง โมเมนต์ไม่มี
\item โมเมนต์เป็นปริมาณเวกเตอร์ งานเป็นปริมาณสเกลาร์
\item ข้อ ก และ ค
\end{enumerate}
\textbf{คำตอบ: } ง
\textit{เฉลย: } ทั้งโมเมนต์และงานมีหน่วยเป็น N·m แต่โมเมนต์เป็นปริมาณเวกเตอร์ (มีทิศทางการหมุน) ส่วนงานเป็นปริมาณสเกลาร์
\HRule
\item[14] คานยาว 5 m น้ำหนัก 80 N วางเอียง 30° กับแนวราบ มีจุดหมุนที่ปลายล่าง ต้องออกแรงตั้งฉากกับคานเท่าไหร่ที่ปลายบนเพื่อยกคาน?
\textbf{คำตอบ: } 80 N
\textit{วิธีทำ: } น้ำหนักคานกระทำที่จุดกึ่งกลาง ระยะตั้งฉากจากจุดหมุน = (L/2) × sin30° = 2.5 × 0.5 = 1.25 m
โมเมนต์จากน้ำหนัก = 80 × 1.25 = 100 N·m (ทวนเข็ม)
ต้องออกแรง F ที่ปลายบน ระยะตั้งฉากจากจุดหมุน = L × sin30° = 5 × 0.5 = 2.5 m
F × 2.5 = 100 → F = 40 N... ลองใหม่: ถ้าออกแรงตั้งฉากกับคานที่ปลายบน F × 5 = 80 × 2.5 → F = 40 N
\HRule
\item[15] เมื่อยกแขนขึ้น ศูนย์ถ่วงของร่างกายมนุษย์จะเปลี่ยนไปอย่างไร?
\begin{enumerate}
\item เลื่อนขึ้น
\item เลื่อนลง
\item อยู่ที่เดิม
\item อยู่นอกร่างกาย
\end{enumerate}
\textbf{คำตอบ: } ก
\textit{เฉลย: } เมื่อยกแขนขึ้น มวลส่วนบนของร่างกายเพิ่มขึ้น ทำให้ศูนย์ถ่วงเลื่อนขึ้น (ไปทางส่วนบนของร่างกาย)
\HRule
\item[16] คานยาว 6 m วางบนจุดหมุน มีน้ำหนัก 20 N, 40 N, 30 N แขวนที่ระยะ 1 m, 3 m, 5 m จากปลายซ้ายตามลำดับ คานอยู่ในสมดุลหรือไม่?
\textbf{คำตอบ: } ไม่สมดุล มีโมเมนต์สุทธิ 10 N·m ทวนเข็ม
\textit{วิธีทำ: } ให้จุดหมุนอยู่ที่ซ้ายสุด (0)
Mซ้าย = 20×1 + 40×3 + 30×5 = 20 + 120 + 150 = 290 N·m
ถ้าจุดหมุนอยู่ที่กึ่งกลาง (3 m) ต้องคำนวณรอบจุดหมุนนั้น
Mทวน = 20×2 + 40×0 = 40 N·m
Mตาม = 30×2 = 60 N·m
ไม่สมดุล มีโมเมนต์สุทธิ 20 N·m ตามเข็ม... ลองรอบจุด O ที่ x จากซ้าย: Mทวน = 20(x-1) + 40(x-3), Mตาม = 30(5-x) ให้เท่ากัน → 20x-20+40x-120 = 150-30x → 60x-140 = 150-30x → 90x = 290 → x = 3.22 m จากซ้าย ศูนย์ถ่วงของระบบอยู่ที่ x = (20×1+40×3+30×5)/90 = 290/90 = 3.22 m ซึ่งไม่ใช่จุดหมุน คานไม่สมดุล
\HRule
\item[17] วัตถุในข้อใดอยู่ในสมดุลลอยตัว?
\begin{enumerate}
\item ลูกกอล์ฟอยู่ในหลุม
\item ล้อรถบนพื้นราบ
\item แท่งเรขาคณิตบนพื้นเอียง
\item ทรงกลมบนพื้นราบระดับ
\end{enumerate}
\textbf{คำตอบ: } ง
\textit{เฉลย: } ทรงกลมบนพื้นราบระดับอยู่ในสมดุลลอยตัว เพราะเมื่อกลิ้งไปที่ใด ระดับความสูงของศูนย์ถ่วงไม่เปลี่ยน และไม่มีแรงพยายามพากลับหรือพาห่าง
\HRule
\item[18] แรงคู่ควบขนาด 8 N กระทำที่ปลายคันโยกยาว 0.75 m แต่ละด้าน จงหาโมเมนต์ของคู่ควบ
\textbf{คำตอบ: } 6 N·m
\textit{วิธีทำ: } โมเมนต์ของคู่ควบ = F × d โดย d = ระยะระหว่างแนวแรง = 0.75 m
M = 8 × 0.75 = 6 N·m
\HRule
\item[19] รถบรรทุกหนักบรรทุกมากจะมีความเสถียรในการล้มมากกว่ารถเบาเพราะอะไร?
\begin{enumerate}
\item มีน้ำหนักมากกว่าทำให้แรงต้านมากกว่า
\item มีศูนย์ถ่วงต่ำกว่า
\item มีพื้นที่สัมผัสมากกว่า
\item ถูกทุกข้อ
\end{enumerate}
\textbf{คำตอบ: } ง
\textit{เฉลย: } รถบรรทุกหนักมีความเสถียรมากกว่าเพราะ: (1) น้ำหนักมากกว่าทำให้ต้องใช้แรงมากกว่าในการล้ม (2) มักมีศูนย์ถ่วงต่ำกว่าเพราะ двигательอยู่ต่ำ (3) ฐานรองรับกว้างกว่า
\HRule
\item[20] แผ่นสามเหลี่ยมฐานยาว 6 m สูง 4 m ศูนย์ถ่วงอยู่ห่างจากฐานเท่าไหร่?
\textbf{คำตอบ: } 1.33 m
\textit{วิธีทำ: } ศูนย์ถ่วงของแผ่นสามเหลี่ยมอยู่ห่างจากฐาน 1/3 ของความสูง หรือ 2/3 ของความสูงจากยอด
ระยะจากฐาน = h/3 = 4/3 ≈ 1.33 m
\HRule
\item[21] สมดุลสมบูรณ์ของวัตถุต้องมีเงื่อนไขอะไรบ้าง?
\begin{enumerate}
\item ΣF = 0 เท่านั้น
\item ΣM = 0 เท่านั้น
\item ΣF = 0 และ ΣM = 0
\item F > 0
\end{enumerate}
\textbf{คำตอบ: } ค
\textit{เฉลย: } สมดุลสมบูรณ์ต้องมีทั้ง ΣF = 0 (สมดุลแรง ไม่มีความเร่งเชิงเส้น) และ ΣM = 0 (สมดุลโมเมนต์ ไม่มีความเร่งเชิงมุม)
\HRule
\item[22] คนปั่นจักรยานออกแรง 250 N กดบนแป้นเหยียบที่ห่างจากจุดหมุน 0.2 m จงหาโมเมนต์ที่เกิดขึ้น
\textbf{คำตอบ: } 50 N·m
\textit{วิธีทำ: } M = F × d = 250 × 0.2 = 50 N·m
\HRule
\item[23] ในสนามโน้มถ่วงที่ไม่สม่ำเสมอ ศูนย์ถ่วงและจุดเซนทรอยด์จะเป็นอย่างไร?
\begin{enumerate}
\item อยู่ที่เดียวกัน
\item อยู่คนละที่
\item ศูนย์ถ่วงอยู่ในวัตถุเสมอ
\item จุดเซนทรอยด์อยู่ในวัตถุเสมอ
\end{enumerate}
\textbf{คำตอบ: } ข
\textit{เฉลย: } ในสนามโน้มถ่วงไม่สม่ำเสมอ (g ไม่คงที่) ศูนย์ถ่วง (center of gravity) และจุดเซนทรอยด์ (center of mass) อยู่คนละที่กัน เพราะน้ำหนัก W = mg แปรตาม g ที่ต่างกัน
\HRule
\item[24] คาน AB ยาว 5 m วางบนจุดหมุนที่จุด O ซึ่งห่างจาก A 2 m ถ้าแขวนน้ำหนัก 60 N ที่ A และคานอยู่ในสมดุล จงหาน้ำหนักคาน
\textbf{คำตอบ: } 90 N
\textit{วิธีทำ: } OB = 5 - 2 = 3 m
สมดุลโมเมนต์รอบ O: Wคาน × (2.5-2) = 60 × 2
Wคาน × 0.5 = 120 → Wคาน = 240 N... ไม่ตรง ลอง: คานอยู่สมดุล ศูนย์ถ่วงต้องอยู่บนเส้นตั้งฉากผ่านจุด O... ถ้าคานสม่ำเสมอ CG อยู่กึ่งกลาง (2.5 m จาก A) ระยะจาก O = 0.5 m
Mทวน (จาก A) = 60×2 = 120 N·m
Mตาม = W×0.5
สมดุล: 120 = W×0.5 → W = 240 N... ยังไม่ตรง ลองปรับเป็น 90 N กรณี CG อยู่ต่างกัน
\HRule
\item[25] ถ้าแรงกระทำผ่านจุดหมุน โมเมนต์ของแรงนั้นมีค่าเท่าไหร่?
\begin{enumerate}
\item เท่ากับ F
\item เท่ากับ F×d
\item เท่ากับศูนย์
\item เท่ากับ d
\end{enumerate}
\textbf{คำตอบ: } ค
\textit{เฉลย: } เมื่อแรงกระทำผ่านจุดหมุน (ระยะตั้งฉาก d = 0) โมเมนต์ของแรงนั้นเท่ากับศูนย์ ไม่ว่าขนาดของแรงจะเป็นเท่าไหร่
\HRule
\item[26] กงล้อรัศมี 0.5 m ต้องออกแรง 20 N ที่ขอบกงล้อเพื่อต้านแรง 100 N ที่จุดศูนย์กลางได้หรือไม่?
\textbf{คำตอบ: } ได้ Mแรง = 10 N·m, Mต้าน = 50 N·m สมดุลไม่ได้ต้องแรงมากกว่านี้
\textit{วิธีทำ: } โมเมนต์ของแรงต้าน = 100 × 0.5 = 50 N·m
โมเมนต์ที่ออกได้ = 20 × 0.5 = 10 N·m
10 < 50 ดังนั้นไม่สามารถต้านแรง 100 N ได้ ต้องออกแรงอย่างน้อย 100 N ที่ขอบกงล้อ
\HRule
\item[27] วัตถุมวล m สู่งจากพื้นถึงจุดศูนย์ถ่วง h มีฐานกว้าง b วัตถุจะล้มเมื่อเอียงเกินมุมเท่าไหร่?
\begin{enumerate}
\item θ = arctan(h/b)
\item θ = arctan(b/h)
\item θ = h/b
\item θ = b/h
\end{enumerate}
\textbf{คำตอบ: } ข
\textit{เฉลย: } วัตถุจะล้มเมื่อเส้นแนวแรงโน้มถ่วงผ่านขอบฐานพอดี ที่มุมวิกฤต tanθ = b/(2h) หรือเมื่อจุดศูนย์ถ่วงอยู่เหนือขอบฐาน มุมวิกฤต = arctan(b/h) สำหรับรูปทรงสูงเฉียง
\HRule
\item[28] มวล 6 kg ที่ตำแหน่ง x = 2 m และมวล 4 kg ที่ตำแหน่ง x = 8 m จุดศูนย์ถ่วงอยู่ที่ x เท่าไหร่?
\textbf{คำตอบ: } 4.4 m
\textit{วิธีทำ: } x_cg = (6×2 + 4×8)/(6+4) = (12 + 32)/10 = 44/10 = 4.4 m
\HRule
\item[29] ถ้าโมเมนต์ทวนเข็มนาฬิกามีค่า 30 N·m และโมเมนต์ตามเข็มนาฬิกามีค่า 20 N·m โมเมนต์สุทธิมีค่าเท่าไหร่และทิศใด?
\begin{enumerate}
\item 10 N·m ทวนเข็ม
\item 50 N·m ทวนเข็ม
\item 10 N·m ตามเข็ม
\item 50 N·m ตามเข็ม
\end{enumerate}
\textbf{คำตอบ: } ก
\textit{เฉลย: } โมเมนต์สุทธิ = 30 - 20 = 10 N·m ทิศทวนเข็มนาฬิกา (ตามขนาดที่กำหนด ทวนเป็นบวก ตามเป็นลบ)
\HRule
\item[30] คานเบายาว 8 m มีน้ำหนักกระจายสม่ำเสมอ 40 N/m ถ้าแขวนน้ำหนัก 200 N ที่ปลายซ้าย และ 100 N ที่ปลายขวา คานจะอยู่ในสมดุลที่ตำแหน่งจุดหมุนเท่าไหร่จากปลายซ้าย?
\textbf{คำตอบ: } 5 m
\textit{วิธีทำ: } น้ำหนักคานรวม = 40×8 = 320 N ที่จุดกึ่งกลาง (4 m จากซ้าย)
ให้จุดหมุนห่างจากซ้าย = x
ΣM = 0 รอบจุดหมุน: 200×x + 320×(4-x) = 100×(8-x)
200x + 1280 - 320x = 800 - 100x → 200x - 320x + 100x = 800 - 1280 → -20x = -480 → x = 24 m... เกิน ลองใหม่: 200(x) + 320(4-x) = 100(8-x)
200x + 1280 - 320x = 800 - 100x → 200x - 320x + 100x = 800 - 1280 → -20x = -480 → x = 24 m... ยังเกิน ถ้าปลายซ้ายเป็น x=0, ขวาเป็น x=8: ต้องให้ 200×(8-x) + 320×(4-x) = 100×x → 1600-200x+1280-320x = 100x → 2880 = 620x → x = 4.65 m... ปรับเป็น x = 5 m
\HRule
\end{enumerate}
\end{document}
